\documentclass[11pt,a4paper]{article}
\usepackage[T1]{fontenc}
\usepackage[utf8]{inputenc}
\usepackage{lmodern}
\usepackage[margin=27mm]{geometry}
\usepackage{amsmath,amssymb,amsthm}
\usepackage{microtype}
\usepackage[hidelinks]{hyperref}
\newtheorem{theorem}{Theorem}[section]
\newtheorem{proposition}[theorem]{Proposition}
\title{Classes and Combination of Upper Bounds for\protect\\
Two- and Three-Dimensional Knapsack Packing}
\author{}
\date{}
\begin{document}
\maketitle
\begin{abstract}
We describe a unified family of geometric relaxations for orthogonal
knapsack packing with permitted rotations and rectangular containers.
The formulation combines volume capacity, transformed-volume inequalities,
data-dependent scale tightening, grouped dissection inequalities, and
orientation incompatibilities. We give validity proofs, distinguish
integer geometric relaxations from their continuous relaxations, and
establish the dominance relations obtained by combining their constraints.
A final extension transfers the per-bin inequalities to bin-packing
lower bounds. The presentation specifies the implemented constructions
without asserting novelty for the underlying dual-feasible functions.
\end{abstract}
% Reusable manuscript section. See main.tex for packages and theorem declarations.
\section{Upper bounds for orthogonal knapsack packing with rotations}
\label{sec:kp-bounds}

\subsection{Problem definition and notation}
Let $d\in\{2,3\}$ and let $C=(C_1,\ldots,C_d)$ be the dimensions of a
rectangular container, with $V=\prod_{j=1}^d C_j$. Each item $i\in I$ has
positive dimensions $a_i=(a_{i1},\ldots,a_{id})$, volume
$v_i=\prod_j a_{ij}$, and profit $p_i\geq0$. Multiplicities are represented
by distinct copies. A feasible solution selects items and places them
axis-parallel inside the container, with disjoint interiors. Let
$\mathcal O_i$ contain the distinct permitted axis permutations that fit
individually into $C$; items with $\mathcal O_i=\varnothing$ are removed.
Write $a_{ioj}$ for the extent of item $i$ along axis $j$ in orientation
$o\in\mathcal O_i$. Thus the formulation covers fixed orientations,
restricted rotations, and all orthogonal rotations. No equality between
the container dimensions is assumed.

The binary variable $y_{io}$ selects item $i$ in orientation $o$, and
\begin{equation}
 x_i=\sum_{o\in\mathcal O_i}y_{io},\qquad
 \sum_{o\in\mathcal O_i}y_{io}\leq1.                 \label{eq:orientation}
\end{equation}
The objective is $\sum_i p_i x_i$. Denote the optimal geometric packing
profit by $z^\star$. All models below omit placement coordinates and
impose necessary conditions for packing. Consequently their maximization
optima are \emph{upper} bounds on $z^\star$, even when their variables
remain binary. The term \emph{integer geometric relaxation} is used for
this distinction; its LP relaxation additionally drops integrality.

\subsection{Area or volume bounds}
The basic necessary condition is
\begin{equation}
 \sum_i v_i x_i\leq V.                             \label{eq:volume}
\end{equation}
Maximizing profit under \eqref{eq:volume} and $x_i\in\{0,1\}$ defines
$U_{\mathrm{vol}}^{\mathrm{IP}}$, an ordinary one-dimensional knapsack
bound. Replacing the binary domain by $0\leq x_i\leq1$ defines the
continuous bound $U_{\mathrm{vol}}^{\mathrm{LP}}$.
In two dimensions ``volume'' means area. Both versions retain the
original item volumes, irrespective of any subsequent dissection.

\subsection{Conservative scales and transformed-volume inequalities}
\label{subsec:scales}
For axis $j$, define the full candidate length domain
\begin{equation}
 S_j=\{a_{ioj}:i\in I,\ o\in\mathcal O_i\}.
\end{equation}
A nonnegative map $f_j:S_j\to\mathbb R_+$ with positive capacity $B_j$
is admissible if
\begin{equation}
 \sum_{s\in S_j}s n_s\leq C_j
 \quad\Longrightarrow\quad
 \sum_{s\in S_j} f_j(s)n_s\leq B_j
 \qquad(n\in\mathbb Z_+^{S_j}).                    \label{eq:admissible}
\end{equation}
The normalized map is $g_j(s)=f_j(s)/B_j$. Requiring the same property
on all lengths in $[0,C_j]$ gives the usual scaled dual-feasible-function
(DFF) condition. Property \eqref{eq:admissible} on the finite candidate
domain is sufficient here. It allows arbitrary repetitions of lengths,
and hence includes every actual pattern of selected items.

For a tuple $t$ of admissible axis maps define
\begin{equation}
 A^t_{io}=\prod_{j=1}^d f^t_j(a_{ioj}),\qquad
 b_t=\prod_{j=1}^d B^t_j,\qquad
 \alpha^t_{io}=\frac{A^t_{io}}{b_t}.
                                                               \label{eq:coefficients}
\end{equation}
Here $b_t$ is the \emph{right-hand-side capacity of row $t$}; it is not
a number of bins. We use $m$ for a bin count below.

\begin{proposition}[Transformed-volume condition]
Every feasible packing satisfies, for every tuple $t$,
\begin{equation}
 \sum_i\sum_{o\in\mathcal O_i} A^t_{io}y_{io}\leq b_t,
 \quad\text{equivalently}\quad
 \sum_i\sum_{o\in\mathcal O_i}\alpha^t_{io}y_{io}\leq1.
                                                               \label{eq:product}
\end{equation}
\end{proposition}
\begin{proof}
Fix a feasible packing and its chosen orientations. For each axis $j$,
put a directed arc from item $i$ to item $h$ whenever the end of $i$'s
projection precedes or equals the start of $h$'s projection. This graph
is acyclic because the original lengths are positive. Every directed
path consists of disjoint intervals and has total original length at
most $C_j$. By \eqref{eq:admissible}, its total transformed length is at
most one when item lengths are replaced by $g_j(a_{ioj})$.

Place each transformed interval at the maximum total transformed length
of a path of predecessors ending immediately before it, or at zero
when there is no predecessor. Its endpoint is at most one. Every arc
still separates the corresponding intervals. Any two original boxes
were separated along at least one axis, so the transformed boxes also
have disjoint interiors. They lie in the unit $d$-cube and therefore
the sum of their volumes is at most one. This is \eqref{eq:product}.
Zero transformed lengths cause no difficulty, since the corresponding
boxes have zero volume.
\end{proof}
This argument is an instance of the conservative-scale framework of
Fekete and Schepers~\cite{FeketeSchepers2004}. It applies after fixing
the orientations in any feasible solution; consequently the same rows
are valid when orientations are decision variables.

\paragraph{Threshold maps.}
For $0\leq q\leq C/2$, use capacity $B=C$ and
\begin{equation}
 T_{C,q}(s)=
 \begin{cases}
 0,&0\leq s<q,\\
 s,&q\leq s\leq C-q,\\
 C,&C-q<s\leq C.
 \end{cases}                                                   \label{eq:threshold}
\end{equation}
If a feasible one-dimensional pattern contains a length exceeding
$C-q$, every other length is less than $q$ and receives zero weight.
Otherwise each transformed length is at most its original value.
Thus the map is admissible. The case $q=0$ is the identity and recovers
\eqref{eq:volume}. For integer dimensions the implementation enumerates
$q\in\{0,\tfrac12,1,\ldots,C/2\}$ and removes duplicate maps.

\paragraph{Floor maps.}
For integer $C\geq1$ and $k\in\{1,\ldots,C\}$, put
$K=\lfloor C/k\rfloor$ and use
\begin{equation}
 F_{C,k}(s)=\left\lfloor s/k\right\rfloor,
 \qquad B=K.                                                   \label{eq:floor}
\end{equation}
Admissibility follows from
$\sum_\ell\lfloor s_\ell/k\rfloor\leq
\lfloor\sum_\ell s_\ell/k\rfloor\leq K$.

\paragraph{Strengthened rounding maps.}
With the same $C,k,K$, use capacity $B=2K$ and
\begin{equation}
 R_{C,k}(s)=
 \begin{cases}
 2\lfloor s/k\rfloor,&2s<C,\\
 K,&2s=C,\\
 2\bigl(K-\lfloor(C-s)/k\rfloor\bigr),&2s>C.
 \end{cases}                                                   \label{eq:rounding}
\end{equation}
\begin{proposition}
The map in \eqref{eq:rounding} is admissible. Moreover,
$R_{C,k}(s)/(2K)\geq F_{C,k}(s)/K$ for every $s\in[0,C]$.
\end{proposition}
\begin{proof}
A pattern contains at most one length $s>C/2$. If it contains one,
the remaining lengths are less than $C/2$, and their total floor
weight is at most $\lfloor(C-s)/k\rfloor$. The total rounding weight
is therefore at most $2K$. Without a length greater than $C/2$,
consider the number of lengths equal to $C/2$. With none, the ordinary
floor-sum inequality gives the result. With one, the other lengths
contribute at most $2\lfloor C/(2k)\rfloor\leq K$. With two, these
lengths contribute exactly $2K$, and only zero lengths can remain.
For dominance, the maps coincide after normalization below $C/2$;
at equality use $2\lfloor C/(2k)\rfloor\leq K$, and above $C/2$ use
$\lfloor s/k\rfloor+\lfloor(C-s)/k\rfloor\leq K$.
\end{proof}
These explicit maps are used as DFF-based ingredients, not asserted
to be new function families. Related strengthened bounds for fixed
orientations are developed by Carlier et al.~\cite{CarlierEtAl2007}.
The formula and proof above specify precisely the variant used here.

\paragraph{Mixed products.}
The axes may use different families and parameters. For example, a
three-dimensional row can use a threshold map on axis 1, a rounding
map on axis 2, and a floor map on axis 3. All coefficients in a row
are evaluated for the \emph{same} orientation variable $y_{io}$.
Likewise all rows share that variable. Allowing each row to choose a
different orientation for an item would discard this consistency
and generally weaken the relaxation.

\subsection{Data-dependent capacity tightening}
For an admissible map $f$ on an axis domain $S$, define
\begin{equation}
 B^*(f;S,C)=\max\left\{\sum_{s\in S}f(s)n_s:
        \sum_{s\in S}s n_s\leq C,\quad n\in\mathbb Z_+^S\right\}.
                                                               \label{eq:tight-capacity}
\end{equation}
For a nonzero map, $0<B^*\leq B$. Substituting $B^*$ for $B$ in
\eqref{eq:admissible} is valid by definition. The product construction
therefore remains valid, and the normalized coefficients can only
increase. For positive integer lengths and capacity, the exact value
is obtained by the unbounded-knapsack recurrence
\begin{align}
 D(0)&=0,\nonumber\\
 D(c)&=\max\left(\{D(c-1)\}\cup
          \{D(c-s)+f(s):s\in S,\ s\leq c\}\right),\quad 1\leq c\leq C,
                                                               \label{eq:dp}\\
 B^*(f;S,C)&=D(C).\nonumber
\end{align}
Unbounded multiplicities deliberately enlarge the possible patterns,
so actual item-copy limits cannot invalidate the resulting capacity.
Under rotations, $S_j$ must include lengths from \emph{all} feasible
orientations, rather than only a preferred orientation. This
construction is a domain-dependent conservative scale; it is not
claimed to be a DFF for arbitrary lengths outside $S_j$, or to exhaust
all data-dependent DFF constructions.

\subsection{Grouped dissection bounds}
\label{subsec:fragments}
Dissect item $i$ into squares for $d=2$ or cubes for $d=3$. Let
$\mathcal Q_i$ be the resulting multiset of side lengths. A complete
dissection satisfies $\sum_{s\in\mathcal Q_i}s^d=v_i$. Any feasible
packing of the original items induces a packing of their fragments;
allowing the fragments to move independently is a relaxation.
Squares and cubes are invariant under orthogonal rotations, so the
same multiset is usable for every orientation of its parent item.

\begin{proposition}[Grouping]
Suppose a fragment inequality assigns nonnegative weight $w_r(s)$
to a fragment of side $s$ and has capacity $\beta_r>0$ in one
container. Then every feasible original-item packing satisfies
\begin{equation}
 \sum_i h_i^r x_i\leq\beta_r,
 \qquad h_i^r=\sum_{s\in\mathcal Q_i}w_r(s).       \label{eq:grouping}
\end{equation}
\end{proposition}
\begin{proof}
Apply the fragment inequality to the packing induced by the selected
original items and collect all terms belonging to each item.
\end{proof}
All sums over $\mathcal Q_i$ count multiplicities. Fragments have no
independent selection variables or profits: every retained fragment
of item $i$ is controlled by $x_i$, and $p_i$ is counted once. Dropping
fragments from a nonnegative row only weakens it. In particular, the
implementation omits unit fragments from these rows while retaining
\eqref{eq:volume} with the full original volumes.

\paragraph{Two-dimensional square bounds.}
The CUTSQ construction repeatedly removes the largest squares that
fit along the longer side of a rectangle and recurses on the residual
rectangle. This is the square-dissection approach used by
Dell'Amico et al.~\cite{DellAmicoEtAl2002}. Let $C=(W,H)$ with $W\geq H$.
The implemented grouped rows include
\begin{align}
 \sum_i x_i\sum_{\substack{s\in\mathcal Q_i\\2s>H}}s&\leq W,
                                                               \label{eq:tall-width}\\
 \sum_i x_i\,\#\{s\in\mathcal Q_i:2s>H\}
    &\leq\left\lfloor\frac{W}{\lfloor H/2\rfloor+1}\right\rfloor,
                                                               \label{eq:tall-count}
\end{align}
where \eqref{eq:tall-count} assumes integer side lengths. Tall squares
must have disjoint horizontal projections. Their widths sum to at
most $W$, and each width is at least $\lfloor H/2\rfloor+1$.
For integer $0\leq q\leq\lfloor H/2\rfloor$, a residual-space row is
\begin{equation}
 \sum_i x_i\left[
    \sum_{\substack{s\in\mathcal Q_i\\s\geq q}}s^2
   +\sum_{\substack{s\in\mathcal Q_i\\s>H-q}}s(H-s)
             \right]\leq WH.                                 \label{eq:residual}
\end{equation}
To verify it, reserve a full-height strip of width $s$ for every
square with $s>H-q$. These strips are mutually disjoint. No other
square of side at least $q$ can overlap such a strip horizontally:
its height plus $s$ exceeds $H$. The strips and the remaining counted
squares therefore have total area at most $WH$, which gives the row.
The benchmark implementation uses the single-bin specialization of
the cited construction under $s\leq W/2$ for all retained squares;
this condition was checked on its inputs. It does not implement every
bound in that paper. The direct arguments above establish validity
of the displayed rows themselves.

\paragraph{Three-dimensional cube bounds.}
One possible dissection sorts a residual box as $a\geq b\geq s$,
sets $k_a=\lfloor a/s\rfloor$, $k_b=\lfloor b/s\rfloor$,
$r_a=a-k_as$, and $r_b=b-k_bs$, and removes $k_ak_b$ cubes of side $s$.
It recurses on the nonempty residual boxes $(r_a,b,s)$ and
$(k_as,r_b,s)$. These boxes are disjoint and conserve volume.
For each axis $j$, the projection inequality is
\begin{equation}
 \sum_i x_i\sum_{\substack{s\in\mathcal Q_i\\
                     2s>C_\ell\ \text{for all }\ell\ne j}}s\leq C_j.
                                                               \label{eq:cube-projection}
\end{equation}
Any two counted cubes overlap in their projections on every other
axis and must therefore be separated along $j$. Also,
\begin{equation}
 \sum_i x_i\,\#\{s\in\mathcal Q_i:2s>\max_j C_j\}\leq1,
                                                               \label{eq:cube-clique}
\end{equation}
since two such cubes cannot be separated along any axis. Finally,
threshold-product rows on the fragments give
\begin{equation}
 \sum_i x_i\sum_{s\in\mathcal Q_i}
           \prod_{j=1}^3 T_{C_j,q_j}(s)\leq V.                  \label{eq:cube-product}
\end{equation}
The same product theorem applies to these cubes. Expanded mixed and
data-tightened families are applied to the original shapes in the
implemented stack; the grouped cube baseline uses threshold products.

\subsection{Incompatibility and clique inequalities}
Two distinct individually fitting items in fixed orientations $(i,o)$
and $(h,r)$ are incompatible precisely when
\begin{equation}
 a_{ioj}+a_{hrj}>C_j\qquad\text{for every }j.                    \label{eq:incompatible}
\end{equation}
Necessity follows because any two nonoverlapping axis-parallel boxes
must be separated along some axis. Conversely, if their lengths sum
to at most $C_j$ on one axis, placing them consecutively on that axis
and at the origin on all other axes packs the pair. Thus every pair
satisfying \eqref{eq:incompatible} yields
\begin{equation}
 y_{io}+y_{hr}\leq1.                                           \label{eq:pair-cut}
\end{equation}
In the item conflict graph, connect $i$ and $h$ only if every pair
of their feasible orientations is incompatible. Every clique $Q$
then yields
\begin{equation}
 \sum_{i\in Q}x_i\leq1.                                       \label{eq:item-clique}
\end{equation}
A subset of cliques suffices for validity. The implementation includes
all incompatible orientation pairs and a greedily generated collection
of maximal item cliques; it does not enumerate all cliques.

\subsection{Combined integer and continuous bounds}
Let $\mathcal T$ be the selected transformed-volume rows,
$\mathcal R$ the grouped fragment rows, $\mathcal E$ the incompatible
orientation pairs, and $\mathcal K$ the selected item cliques. The
combined integer bound is the optimal value of
\begin{subequations}\label{eq:combined}
\begin{align}
 U^{\mathrm{IP}}=\max\quad&\sum_i p_i x_i                       \label{eq:obj}\\
 \text{s.t.}\quad
 &x_i=\sum_{o\in\mathcal O_i}y_{io}\leq1 && i\in I,\\
 &\sum_i v_i x_i\leq V,\\
 &\sum_i\sum_{o\in\mathcal O_i}\alpha^t_{io}y_{io}\leq1
                                      &&t\in\mathcal T,\\
 &\sum_i h_i^r x_i\leq\beta_r &&r\in\mathcal R,\\
 &y_{io}+y_{hr}\leq1 &&((i,o),(h,r))\in\mathcal E,\\
 &\sum_{i\in Q}x_i\leq1 &&Q\in\mathcal K,\\
 &y_{io}\in\{0,1\} &&i\in I,\ o\in\mathcal O_i.
\end{align}
\end{subequations}
The LP bound $U^{\mathrm{LP}}$ replaces the last line by $y_{io}\geq0$;
the item constraints already imply $y_{io}\leq1$. The equality defining
$x_i$ permits its elimination in both models.

\begin{theorem}[Validity and dominance]\label{thm:dominance}
For any finite selection of the valid rows described above,
\begin{equation}
 z^\star\leq U^{\mathrm{IP}}\leq U^{\mathrm{LP}}
                         \leq U_{\mathrm{vol}}^{\mathrm{LP}},
 \qquad U^{\mathrm{IP}}\leq U_{\mathrm{vol}}^{\mathrm{IP}}.
                                                               \label{eq:dominance}
\end{equation}
Adding valid rows cannot increase either bound.
\end{theorem}
\begin{proof}
Every geometric solution induces a binary assignment satisfying all
rows, which proves the first inequality. Relaxing integrality enlarges
the feasible set, proving the second. Projecting any combined LP
solution onto $x$ gives a feasible fractional volume-knapsack solution;
projecting an integer solution gives a feasible binary volume-knapsack
solution. The objective is preserved, proving both volume comparisons.
Adding rows restricts the relevant feasible set.
\end{proof}
There is no corresponding general claim that $U^{\mathrm{LP}}$ is at
most $U_{\mathrm{vol}}^{\mathrm{IP}}$, nor that a shape-only relaxation
dominates every grouped-fragment relaxation. Their intersection is
used to obtain both sets of restrictions.

\paragraph{Combination by intersection and by minimum.}
If each class is solved separately to produce $U_k$, then $\min_k U_k$
is a valid upper bound. Combining their rows in one model gives
\begin{equation}
 U_{\cap}^{\mathrm{IP}}\leq\min_k U_k^{\mathrm{IP}},\qquad
 U_{\cap}^{\mathrm{LP}}\leq\min_k U_k^{\mathrm{LP}}.             \label{eq:intersection}
\end{equation}
The inequalities can be strict because separate classes can admit
different maximizing selections or orientations. Bounds obtained
separately are not added together.

\paragraph{Nested implementation.}
The baseline contains the volume row, original-shape threshold products,
and grouped square or cube rows. The next stages add, in order,
products from the axis families $\{T,F\}$, products from $\{T,F,R\}$,
products of capacity-tightened maps, and incompatibility/clique rows.
Each stage retains all previous constraints. Consequently, for
$\bullet\in\{\mathrm{IP},\mathrm{LP}\}$,
\begin{equation}
 U_0^\bullet\geq U_1^\bullet\geq U_2^\bullet
                      \geq U_3^\bullet\geq U_4^\bullet\geq z^\star.
\end{equation}
A map $f/B$ may be discarded when another map $\widetilde f/\widetilde B$
is componentwise at least as large on the entire axis domain. With
nonnegative factors, every product using the former is dominated by
the corresponding product using the latter. The implementation makes
these comparisons by integer cross multiplication, tightens the
retained axis maps, and removes duplicates. This is a finite family
selection, not optimization over all possible conservative scales.

\paragraph{Certificates and incomplete optimization.}
After eliminating $x$, normalize all geometric rows to write
$Ay\leq\mathbf1$ and retain the item constraints separately. The LP
dual is
\begin{equation}
 \min_{\lambda\geq0,\,\nu\geq0}
 \left\{\sum_t\lambda_t+\sum_i\nu_i:
       \sum_t A_{t,io}\lambda_t+\nu_i\geq p_i
                  \quad\forall i,o\right\}.                  \label{eq:dual}
\end{equation}
Thus nonnegative multipliers combine the row capacities into an
explicit upper-bound certificate. For a time-limited integer solve,
a valid solver global upper bound is usable; its incumbent objective
alone is not an upper bound. An incumbent need not represent a feasible
geometric packing. With integer profits, a mathematically valid real
upper bound $\overline U$ can be replaced by $\lfloor\overline U\rfloor$.
Floating-point feasibility checks and rounding tolerances should be
reported separately from exact mathematical certification.

\subsection{Transfer to bin packing}
\label{subsec:bin-packing}
The same per-container rows also yield lower bounds when every item
must be assigned to one of $m$ identical bins. For an orientation-independent
row $\sum_{i\text{ in one bin}}h_i^r\leq\beta_r$, summing across bins gives
\begin{equation}
 m\geq\left\lceil\frac{\sum_i h_i^r}{\beta_r}\right\rceil.
                                                               \label{eq:bp-simple}
\end{equation}
For orientation-dependent rows, let $A_{t,io}$ denote normalized
coefficients of any selected per-bin inequalities. A consistent
aggregate relaxation is
\begin{equation}
 L^{\mathrm{IP}}=\min\left\{m:
  \sum_{o\in\mathcal O_i}y_{io}=1\ (i\in I),\quad
  \sum_{i,o}A_{t,io}y_{io}\leq m\ (t\in\mathcal H),\quad
  y_{io}\in\{0,1\},\ m\in\mathbb Z_+\right\}.                 \label{eq:bp-model}
\end{equation}
Any $m$-bin packing supplies a feasible assignment by using its actual
orientations and summing each valid row over the bins. Therefore
$L^{\mathrm{IP}}$ is a lower bound on the optimal bin count. Relaxing
$y$ and $m$ to nonnegative continuous variables gives $L^{\mathrm{LP}}$,
and $\lceil L^{\mathrm{LP}}\rceil$ is also a valid lower bound. These
aggregate models impose necessary total capacities; they do not
construct bin assignments. The transfer concerns the underlying
valid inequalities, not an arbitrary conversion of a scalar knapsack
profit bound into a bin count.

\bibliographystyle{plain}
\bibliography{references}
\end{document}
